\section{Silogisme Hipotetis dan Silogisme Disjungtif}

Silogisme menghubungkan beberapa proposisi melalui implikasi atau disjungsi. Berbeda dengan Modus Ponens dan Modus Tollens yang bekerja pada satu implikasi utama, silogisme sering digunakan untuk merangkai beberapa premis \cite{rosen2019}.

\subsection{Silogisme Hipotetis}

Bentuk ini memanfaatkan sifat transitif pada implikasi.

\begin{teorema}[Silogisme Hipotetis]
Jika $p \rightarrow q$ dan $q \rightarrow r$, maka dapat disimpulkan $p \rightarrow r$. Proposisi perantara $q$ tidak muncul pada konklusi.
\end{teorema}

\textbf{Skema argumen Silogisme Hipotetis:}
\[
\begin{array}{c}
p \rightarrow q \\
q \rightarrow r \\
\hline
\therefore p \rightarrow r
\end{array}
\]

\begin{contoh}
\begin{itemize}
    \item \textbf{Premis 1 ($p \rightarrow q$)}: ``Jika saya belajar logika setiap malam, maka saya memahami aljabar proposisi.''
    \item \textbf{Premis 2 ($q \rightarrow r$)}: ``Jika saya memahami aljabar proposisi, maka saya lulus ujian.''
    \item \textbf{Konklusi ($\therefore p \rightarrow r$)}: ``Jika saya belajar logika setiap malam, maka saya lulus ujian.''
\end{itemize}
\end{contoh}

\subsection{Silogisme Disjungtif}

Disjungsi menyatakan bahwa paling sedikit salah satu dari dua proposisi bernilai benar.

\begin{teorema}[Silogisme Disjungtif]
Dari $p \vee q$ dan $\neg p$, dapat disimpulkan $q$. Secara simetris, dari $p \vee q$ dan $\neg q$, dapat disimpulkan $p$.
\end{teorema}

\textbf{Skema argumen Silogisme Disjungtif:}
\[
\begin{array}{ccc}
p \vee q & & p \vee q\\
\neg p   & \text{atau} & \neg q \\
\hline
\therefore q & & \therefore p
\end{array}
\]

\begin{contoh}
\begin{itemize}
    \item \textbf{Premis 1 ($p \vee q$)}: ``Pelaku melarikan diri melalui jendela ($p$) atau melalui plafon ($q$).''
    \item \textbf{Premis 2 ($\neg p$)}: ``Tidak ada bukti pelaku melarikan diri melalui jendela.''
    \item \textbf{Konklusi ($\therefore q$)}: ``Pelaku melarikan diri melalui plafon.''
\end{itemize}
\end{contoh}

Bentuk $[(p \vee q) \land \neg p] \rightarrow q$ dapat diverifikasi sebagai tautologi melalui tabel kebenaran empat baris.
